\cotefile{101}
\scanpage{30}
% L1–2
On peut \unsure{remplacer|expliciter|expliquer}, dans le cas général, le th. de dualité locale par une suite spectrale.

% L3–5
\uline{Lemme 6} $A'\longrightarrow\deleted{B}A$ un hom. d'anneaux, $N$ un $\deleted{B}A$-module, $M$ un $A'$-module. On a une suite spectrale cohomologique
% L6
\[
\operatorname{Ext}_{A'}^{\unsure{*|n|\bullet}}(N,M)\Longleftarrow
\operatorname{Ext}_{\deleted{B}A}^{p}(N,\operatorname{Ext}_{A'}^{q}(\deleted{B}A,M))
\]
% margin note written sideways
\margin{\uline{Cor} 1 Soit $x\in A$ non diviseur de $0$ dans $A$, ni dans $M$. Alors on a $\operatorname{Ext}_{A}^{p}(N,M)\simeq\operatorname{Ext}_{A/x}^{p-1}(N,M/xM)$ pour tt $N$, \uline{Cor 2} \unsure{Avec} une $M$-\unsure{suite}.}

% L7–8
Elle provient de la formule d'associativité
\[
\operatorname{Hom}_{A'}(N,M)=\operatorname{Hom}_{\unsure{B|A}}(N,\operatorname{Hom}_{A'}(B,M))
\]
% L9–10
\diagram{\begin{tikzcd}[ampersand replacement=\&]
{\unsure{\mathcal C|\mathscr C|\mathcal E}}^{\unsure{A'|A}}\arrow[r,dash,"{h_{\deleted{B}A}^{A'}}"]\arrow[rr,bend right=23,swap,"{h_N^{A'}}"]
 \&\mathcal C^{\deleted{B}A}\arrow[r,"{h_N^{\deleted{B}A}}"] \&{\unsure{\mathcal G|\mathscr G}}^{\deleted{B}A}
\end{tikzcd}}
$h_{\deleted{B}A}^{A'}$ transforme injectifs en injectifs

\centerline{- - - - -}

% L11–13
\uline{Prop. 12} \inserted{Soit $A$ un anneau quotient d'un Coh. Mac. $A'$ muni d'un mod. fdt $\Omega_{A'}$.} Il y a une suite spectrale canonique de type cohomologique
% L14
\[
\boxed{(E^{\deleted{\unsure{*}}i}(M))=(H^{\deleted{\unsure{d}}-i}(M))\underset{\unsure{p}}{\Longleftarrow}
 E_2^{pq}(M)=D(\operatorname{Ext}_{A}^{p}(M,\Omega_{A}^{(q)}))}
\]
% slanted
\margin{\uline{\unsure{Module fondamental|Modules fondamentaux}} \uline{pour un} \uline{anneau} de \uline{\unsure{polynomes}}\unsure{....}}

% L15–16
On applique le lemme 6 avec $N$\deleted{\unsure{$=$}} remplacé par $M$, et $M$ remplacé par $\Omega_{A'}$.

% archive sheet number: 29
